%! The Quadratic Formula & the Discriminant — Unit 3, Lesson 3.5 — Homework (Answer Key)
\documentclass[10pt]{article}
\usepackage{algebra2-article}
\usepackage{algebra2-key}   % pulls in -boxes; do NOT also load -boxes
\newcommand{\TallMath}[1]{$\displaystyle #1\rule[-1.4em]{0pt}{3.2em}$}

\begin{document}

\pageheader{Unit 3, Lesson 3.5}{Homework --- Answer Key}
\namedateperiod

\vspace{0.10in}

\begin{notesbox}{Practice}
Show your work. Write in standard form first, read $a,b,c$ with signs, and simplify every answer (use
$a+bi$ when the roots are complex).

\begin{enumerate}[leftmargin=1.9em, itemsep=9pt]
  \item \textbf{Standard form, then $a,b,c$.}
    \begin{enumerate}[label=(\alph*), leftmargin=1.9em, itemsep=4pt]
      \item $x^2=6x-2 \Rightarrow$ \ans{$x^2-6x+2=0$}\quad $a=\ans{$1$}\ b=\ans{$-6$}\ c=\ans{$2$}$
      \item $5x-3=2x^2 \Rightarrow$ \ans{$2x^2-5x+3=0$}\quad $a=\ans{$2$}\ b=\ans{$-5$}\ c=\ans{$3$}$
    \end{enumerate}

  \item \textbf{Discriminant \& classify} (compute $b^2-4ac$; state two real / one repeated / two complex).
    \[ x^2-2x-8=0:\ \ans{$36$; two real} \qquad x^2+x+1=0:\ \ans{$-3$; two complex} \]
    \[ 9x^2-12x+4=0:\ \ans{$0$; one repeated} \qquad x^2+5x+2=0:\ \ans{$17$; two real} \]

  \item \textbf{Solve} (real, rational).
    \begin{enumerate}[label=(\alph*), leftmargin=1.9em, itemsep=4pt]
      \item $x^2-x-6=0 \Rightarrow x=$ \ans{$3,\ -2$}
      \item $2x^2-7x+3=0 \Rightarrow x=$ \ans{$3,\ \dfrac{1}{2}$}
    \end{enumerate}

  \item \textbf{Solve} (real, irrational --- simplest radical form).
    \begin{enumerate}[label=(\alph*), leftmargin=1.9em, itemsep=4pt]
      \item $x^2-6x+7=0 \Rightarrow x=$ \ans{$3\pm\sqrt{2}$}
      \item $x^2+4x+1=0 \Rightarrow x=$ \ans{$-2\pm\sqrt{3}$}
    \end{enumerate}

  \item \textbf{Solve} (complex --- $a+bi$ form).
    \begin{enumerate}[label=(\alph*), leftmargin=1.9em, itemsep=4pt]
      \item $x^2-2x+10=0 \Rightarrow x=$ \ans{$1\pm3i$}
      \item $x^2+16=0 \Rightarrow x=$ \ans{$\pm4i$}
      \item $2x^2+2x+1=0 \Rightarrow x=$ \ans{$-\dfrac{1}{2}\pm\dfrac{1}{2}i$}
    \end{enumerate}

  \item \textbf{Discriminant $\to$ graph.} State how many times each parabola crosses the $x$-axis.
    \[ b^2-4ac=49:\ \ans{$2$} \qquad b^2-4ac=0:\ \ans{$1$} \qquad b^2-4ac=-4:\ \ans{$0$} \]

  \item \textbf{Choose the most efficient method} (square roots / factoring / quadratic formula --- you need not solve).
    \[ x^2-25=0:\ \ans{square roots} \quad x^2-7x+12=0:\ \ans{factoring} \quad x^2+3x+5=0:\ \ans{formula} \]

  \item \textbf{Explain.} A classmate says ``the discriminant is negative, so the equation has \emph{no}
        solutions.'' What is the correct statement?
        \ansline{A negative discriminant means no \emph{real} solutions --- but there are still \emph{two} solutions: a complex-conjugate pair $a\pm bi$. Every quadratic has solutions over the complex numbers.}
\end{enumerate}
\end{notesbox}

\begin{extensionbox}
\textbf{1. Model.} A firework is launched so its height is $h(t)=-16t^2+32t+48$ (feet, $t$ in seconds).
Solve $h(t)=0$ with the formula to find when it lands; reject the impossible root.
\[ t=\ans{$\dfrac{-32\pm64}{-32}=3,\ -1$} \Rightarrow t=\ans{$3$}\text{ s} \]

\textbf{2. Parameter.} For what value of $k$ does $x^2+6x+k=0$ have exactly one repeated real root?
\ans{$k=9$}\; For what values of $k$ does it have two complex roots? \ans{$k>9$}

\textbf{3. Verify a complex root.} Substitute $x=1+3i$ into $x^2-2x+10$ (from \#5a) and show it gives $0$.
\ansline{$(1+3i)^2-2(1+3i)+10=(-8+6i)-(2+6i)+10=0$. \checkmark}
\end{extensionbox}

\begin{spiralbox}
\textbf{Coming next --- Lesson 3.6: Systems Involving Quadratics.} You can now solve \emph{any} single
quadratic. Next you'll solve \emph{systems} --- a line and a parabola, or two parabolas --- by substitution.
The catch: substitution collapses the system into one quadratic equation, which you'll finish with today's
formula, and the discriminant will tell you how many times the two graphs intersect.
\end{spiralbox}

\begin{teachernote}
Answer traps to flag while grading: \#1 sign flips moving terms across ($5x-3=2x^2\to2x^2-5x+3=0$); \#2/\#6
the $-4ac$ sign when $c<0$; \#3b/\#5c forgetting $2a$ in the denominator or not reducing; \#4 leaving
$\frac{6\pm2\sqrt2}{2}$ unsimplified instead of $3\pm\sqrt2$; \#5 stopping at $\sqrt{-36}$ without the $i$;
\#8 the ``no solutions'' misconception (it's no \emph{real} solutions). The extension \#1 is a clean $t=3$~s
(the same firework from Lesson~3.0); if a student applies the formula to $-16t^2+32t+48=0$ directly they get
$t=\frac{-32\pm64}{-32}$, and clearing to $t^2-2t-3=0$ factors to $(t-3)(t+1)$ --- a nice
formula-vs-factoring check. \#2 makes the discriminant a dial; \#3 closes back to Lesson~3.4.
\end{teachernote}

\end{document}
